\chapter{Segurança da Informação}
\label{cap:seguranca-teoria}

\begin{edital}
Políticas e procedimentos de segurança; ISO/IEC 27001:2022 e 27002:2022;
confidencialidade, integridade e disponibilidade; mecanismos de segurança
(controle de acesso, OAuth2, SSO); gestão de riscos (ameaça, vulnerabilidade,
impacto); SDL e OWASP Top 10; análise estática e dinâmica de código (SAST e
DAST). Também: HTTPS, SSL/TLS.
\end{edital}

\section{A tríade CIA: o critério que organiza tudo}

\begin{conceito}[Cada controle existe para sustentar um dos três pilares]
A tríade não é uma lista decorativa de três palavras --- é o
\textbf{critério de classificação} de todo o resto do capítulo. Cada
controle (cifra, hash, backup, RBAC, WAF) existe para sustentar um dos
três pilares, e a pergunta que resolve metade das questões do bloco é
sempre a mesma: \textit{qual pilar este controle protege?}

\begin{itemize}
\item \textbf{Confidencialidade}: só quem tem direito acessa/\textbf{vê}
o dado (cifra, controle de acesso, mascaramento).
\item \textbf{Integridade}: o dado permanece \textbf{íntegro}, não
alterado indevidamente (hash, assinatura digital, trilha de auditoria).
\item \textbf{Disponibilidade}: o dado/serviço fica \textbf{acessível}
quando necessário (redundância, backup, plano de continuidade).
\item Complementos frequentes: autenticidade, não repúdio, legalidade.
\end{itemize}

O detalhe que separa quem entendeu de quem decorou: os três pilares
\textbf{competem entre si}. Aumentar confidencialidade quase sempre custa
disponibilidade --- um segundo fator de autenticação barra o invasor e
também atrasa o usuário legítimo que perdeu o celular; cifrar o backup
protege o dado roubado, mas o torna inútil se a chave de cifra se perder.
Por isso uma questão de cenário raramente pede "o mais seguro" --- pede o
\textbf{equilíbrio adequado ao contexto}, e a alternativa que maximiza um
pilar ignorando os outros costuma ser a errada.
\end{conceito}

\begin{cuidado}
Mapeie o verbo do cenário para o pilar certo: "restringir quem vê" é
confidencialidade, não integridade; "ficar no ar" é disponibilidade;
"alteração indevida do conteúdo" é integridade.
\end{cuidado}

\section{Criptografia}

\begin{conceito}[Assimétrica: a mesma matemática, duas direções, dois objetivos]
O par de chaves tem uma propriedade central, e ela é a origem de toda a
confusão sobre o tema: \textbf{o que uma chave fecha, só a outra abre}.
Como as duas direções funcionam, a mesma matemática serve a dois objetivos
diferentes --- e é a \textbf{direção escolhida} que decide qual.

\begin{center}
\begin{tabular}{@{}p{2.6cm}p{4.4cm}p{4.6cm}@{}}
\toprule
& \textbf{Para SIGILO} & \textbf{Para ASSINATURA} \\
\midrule
Cifra com & a \textbf{pública do destinatário} & a \textbf{privada do
emissor} \\
Abre com & a \textbf{privada do destinatário} & a \textbf{pública do
emissor} \\
Quem consegue ler & só o destinatário --- só ele tem a privada &
\textbf{qualquer um}: a pública é pública \\
Logo, garante & \textbf{confidencialidade} & \textbf{autenticidade,
integridade e não repúdio} \\
\bottomrule
\end{tabular}
\end{center}

A lógica de cada uma, resumida numa frase: para \textbf{sigilo}, eu quero
que \textit{só você} leia, então uso o cadeado que \textit{só você} abre
--- sua chave pública. Para \textbf{assinatura}, eu quero provar que
\textit{fui eu}, então uso a chave que \textit{só eu} tenho --- minha
privada; que qualquer um consiga abrir e verificar não é defeito, é o
ponto central: é assim que qualquer pessoa consegue confirmar a
assinatura.

Duas notas importantes. Primeira: assina-se o \textbf{hash} do documento,
não o documento inteiro --- criptografia assimétrica é lenta, e o hash tem
tamanho fixo e pequeno. Segunda: na prática, nem o \textbf{sigilo} usa
assimétrica para o volume de dados. O TLS faz o que se chama de
\textbf{envelope}: usa a assimétrica \textbf{só para combinar} uma chave
simétrica de sessão entre as partes, e daí em diante cifra tudo com a
simétrica, que é rápida. "A assimétrica substituiu a simétrica" é falso
--- elas se combinam, cada uma no papel em que é boa.
\end{conceito}

\begin{conceito}[Simétrica x assimétrica, lado a lado]
\begin{center}
\begin{tabular}{@{}p{2.6cm}p{4.6cm}p{4.6cm}@{}}
\toprule
& \textbf{Simétrica} & \textbf{Assimétrica (chave pública)} \\
\midrule
Chaves & uma chave secreta compartilhada & par: pública + privada \\
Velocidade & rápida & lenta \\
Uso & cifrar o volume de dados & troca de chave, assinatura \\
Exemplos & AES, DES/3DES & RSA, ECC \\
\bottomrule
\end{tabular}
\end{center}
\end{conceito}

\begin{conceito}[Hash, assinatura, certificado, TDE, esteganografia]
\textbf{Hash} (MD5, SHA-256): resumo de tamanho fixo, \textbf{unidirecional}
--- não existe "descriptografar" um hash. Garante \textbf{integridade}, não
confidencialidade (MD5 e SHA-1 hoje são considerados fracos, com colisões
conhecidas). \textbf{Assinatura digital}: o hash do documento é cifrado
com a \textbf{chave privada} do emissor $\to$ garante autenticidade,
integridade e não repúdio ao mesmo tempo. \textbf{Certificado digital /
ICP-Brasil}: vincula uma identidade a uma chave pública, emitido por uma
AC (Autoridade Certificadora) --- é o que resolve o problema de "como sei
que esta chave pública é realmente da pessoa X". \textbf{TDE
(\textit{Transparent Data Encryption})}: cifra os dados \textbf{em
repouso} (arquivos do banco/disco), de forma transparente à aplicação ---
protege contra o \textbf{roubo do arquivo físico}, não contra acesso
lógico já autorizado ao banco. \textbf{Esteganografia} é diferente de
criptografia: ela \textbf{oculta a existência} da mensagem (esconde dado
dentro de uma imagem/áudio), enquanto a criptografia oculta o
\textbf{conteúdo}, sem esconder que há uma mensagem.
\end{conceito}

\begin{conceito}[HMAC: o cenário do \textit{webhook}]
\textbf{HMAC} (\textit{Hash-based Message Authentication Code}) é um hash
\textbf{combinado a uma chave secreta compartilhada}. Como só quem tem a
chave produz o código correto, o HMAC entrega \textbf{integridade} (a
mensagem não foi alterada) \textbf{e} \textbf{autenticidade} (veio de
quem tem a chave). O que ele \textbf{não} entrega: \textbf{confidencialidade}
(a mensagem viaja legível, só o código de autenticação vai junto) e
\textbf{não repúdio} (a chave é \textbf{simétrica}, e as duas pontas
conseguem gerar o mesmo código --- então não dá para provar, perante um
terceiro, qual das duas partes gerou a mensagem).

O cenário clássico: um sistema externo envia um \textit{webhook} (uma
chamada HTTP de notificação) e é preciso garantir que a requisição veio
mesmo do parceiro combinado, e não foi adulterada no caminho. A resposta é
assinar o corpo com \textbf{HMAC} usando um segredo combinado entre os
dois lados, e o receptor recalcular o código e comparar. Hash puro não
serve (qualquer um recalcularia sem precisar de segredo); cifrar o corpo
resolveria confidencialidade, que não é o problema descrito.
\end{conceito}

\begin{cuidado}
HMAC \textbf{não} garante confidencialidade (não há cifra do conteúdo em
si) nem não repúdio (para isso é preciso \textbf{assinatura digital} com
chave \textbf{privada}, que só o emissor possui). Guarde o corte: chave
\textbf{simétrica} $\to$ autenticidade + integridade; chave
\textbf{privada} $\to$ acrescenta o não repúdio.
\end{cuidado}

\section{HTTPS, SSL e TLS}

\begin{conceito}
\textbf{HTTPS = HTTP sobre TLS/SSL}. Garante confidencialidade,
integridade e autenticação do servidor (o cliente sabe que está falando
com o servidor legítimo). \textbf{TLS substitui o SSL}: corrige
vulnerabilidades conhecidas das versões antigas do SSL. SSL está
obsoleto; TLS 1.2/1.3 é o padrão atual.
\end{conceito}

\begin{cuidado}
SSL \textbf{não} é mais seguro/moderno que TLS --- é o contrário: TLS
sucedeu e corrigiu o SSL. Os dois não são intercambiáveis apesar de
frequentemente citados juntos.
\end{cuidado}

\section{Controle de acesso}

\begin{conceito}[Quatro políticas, quatro bases de decisão diferentes]
\begin{center}
\begin{tabular}{@{}p{3cm}p{8.2cm}@{}}
\toprule
\textbf{Política} & \textbf{Base da decisão} \\
\midrule
DAC (discricionário) & o \textbf{dono do recurso} concede acesso, a seu
critério \\
MAC (mandatório) & \textbf{rótulos de segurança} comparados a
autorizações; a regra é central e o usuário \textbf{não decide} \\
RBAC (por papéis) & acesso concedido via \textbf{papéis/funções}
organizacionais \\
ABAC (por atributos) & decisão calculada a partir de \textbf{atributos}
de usuário, recurso e contexto no momento do acesso \\
\bottomrule
\end{tabular}
\end{center}

\textbf{Princípio do menor privilégio}: conceder só o acesso estritamente
necessário para a função (um gerente de RH acessa só o que sua função
exige, não tudo que tecnicamente seria possível liberar).
\end{conceito}

\begin{conceito}[OAuth2, OIDC e SSO]
\textbf{OAuth2}: protocolo de \textbf{autorização delegada} baseado em
tokens --- resolve "este aplicativo pode agir em meu nome, com este
escopo de permissão limitado", \textbf{sem} entregar a senha ao
aplicativo. É diferente de autenticação. \textbf{OpenID Connect (OIDC)}:
construído \textbf{sobre} o OAuth2, é quem efetivamente faz
\textbf{autenticação} --- adiciona um ID Token com \textit{claims}
padronizados: \textbf{iat} (\textit{issued at}, momento da emissão),
\textbf{exp} (expiração), \textbf{sub} (\textit{subject}, identificador
único do usuário), \textbf{jti} (identificador único do próprio token).
\textbf{SSO (\textit{Single Sign-On})}: um único login dá acesso a vários
sistemas, sem repetir autenticação em cada um.
\end{conceito}

\begin{cuidado}
MAC significa rótulos e classificações impostos pelo \textbf{sistema} (o
usuário não decide, diferente do DAC, em que o \textbf{proprietário}
decide). OAuth2 é \textbf{autorização}, não autenticação --- confundir os
dois é o erro mais comum do tema.
\end{cuidado}

\section{OWASP Top 10}

\begin{conceito}[Duas edições vigentes: 2025 (atual) e 2021 (a mais citada em prova)]
A edição \textbf{vigente} é a \textbf{2025} (oitava da série). A
\textbf{2021} continua relevante como referência amplamente citada em
provas. Ao ler qualquer enunciado sobre o tema, \textbf{verificar qual ano
ele cita} antes de contar posição é essencial, porque a ordem e até a
existência de categorias mudou entre as duas edições.

\begin{center}
\begin{tabular}{@{}p{1.3cm}p{5.3cm}p{5.3cm}@{}}
\toprule
\# & \textbf{Top 10:2025} & \textbf{Top 10:2021} \\
\midrule
1 & Broken Access Control (absorveu o SSRF) & Broken Access Control \\
2 & Security Misconfiguration (subiu de 5º) & Cryptographic Failures \\
3 & Software Supply Chain Failures (nova) & Injection (SQLi; \textbf{XSS
absorvido aqui}) \\
4 & Cryptographic Failures (caiu de 2º) & Insecure Design (nova em 2021) \\
5 & Injection (caiu de 3º) & Security Misconfiguration \\
6 & Insecure Design & Vulnerable and Outdated Components \\
7 & Authentication Failures & Identification and Authentication Failures \\
8 & Software \textit{or} Data Integrity Failures & Software \textit{and}
Data Integrity Failures \\
9 & Security Logging and \textbf{Alerting} Failures & Security Logging
and \textbf{Monitoring} Failures \\
10 & Mishandling of Exceptional Conditions (nova) & Server-Side Request
Forgery (SSRF) \\
\bottomrule
\end{tabular}
\end{center}

O que mudou de 2021 para 2025: \textbf{duas categorias novas} (Software
Supply Chain Failures e Mishandling of Exceptional Conditions) e
\textbf{uma consolidação} (o SSRF, que era categoria própria em 2021, foi
absorvido pelo Broken Access Control). Três categorias só \textbf{mudaram
de nome}, mantendo a essência.
\end{conceito}

\begin{cuidado}
Desde 2021, \textbf{XSS entrou dentro de Injection} (era categoria própria
até 2017) --- não é mais listado separadamente. Itens como "proteção da
cadeia de suprimentos/pipeline", "treinamento operacional" pertencem a
outros frameworks (CI/CD, SSDF), não ao OWASP Top 10 do lado
\textit{web} tradicional --- exceto que, a partir de 2025, cadeia de
suprimentos \textbf{de software} virou categoria própria (A03), o que não
retroage para tornar isso item de 2021.
\end{cuidado}

\section{Desenvolvimento seguro}

\begin{conceito}
\textbf{SDL} (\textit{Security Development Lifecycle}): incorpora
segurança em \textbf{todo} o ciclo de desenvolvimento, não só numa
revisão final. \textbf{SAST} (estática): analisa o \textbf{código-fonte}
\textbf{sem executá-lo} --- é caixa-branca. \textbf{DAST} (dinâmica):
testa a \textbf{aplicação em execução}, sem olhar o código-fonte --- é
caixa-preta. \textbf{IAST}: combina os dois, instrumentando a aplicação
em tempo de execução para observar o comportamento interno enquanto ela
roda. \textbf{DevSecOps}: automatiza essas verificações de segurança
dentro do próprio pipeline de CI/CD.
\end{conceito}

\begin{cuidado}
SAST analisa código \textbf{parado}; DAST testa aplicação \textbf{rodando}
--- não inverter os dois.
\end{cuidado}

\section{X.800: a arquitetura de segurança OSI}

\begin{conceito}[Vocabulário comum a normas e bancas]
A recomendação X.800 (ITU-T) dá o vocabulário de segurança usado por
diversas normas posteriores. Divide o assunto em três: \textbf{ataques},
\textbf{serviços} e \textbf{mecanismos}. O foco cobrado em prova é a
divisão dos \textbf{mecanismos}.

Os cinco \textbf{serviços} de segurança: autenticação, controle de
acesso, confidencialidade, integridade e \textbf{não repúdio} (a
disponibilidade aparece como categoria adicional, fora dos cinco
clássicos).

\textbf{Mecanismos}, divididos em dois grupos:

\begin{center}
\begin{tabular}{@{}p{3.4cm}p{8.6cm}@{}}
\toprule
\textbf{Específicos} (ligados a uma camada e a um serviço) & cifração,
assinatura digital, controle de acesso, integridade de dados, troca de
autenticação, \textbf{preenchimento de tráfego} (\textit{traffic
padding}), controle de roteamento e \textbf{notarização} \\
\textbf{Disseminados} (\textit{pervasive}, não específicos de uma camada)
& funcionalidade confiável, \textbf{rótulos de segurança}, detecção de
eventos, \textbf{trilha de auditoria} de segurança e recuperação de
segurança \\
\bottomrule
\end{tabular}
\end{center}

O critério de corte: um mecanismo \textbf{específico} implementa
\textbf{um serviço} numa \textbf{camada} determinada; um mecanismo
\textbf{disseminado} é de gestão, e vale para o sistema inteiro, sem se
prender a uma camada só. \textbf{Preenchimento de tráfego} e
\textbf{controle de roteamento} costumam surpreender por serem
específicos --- eles servem à confidencialidade do \textbf{fluxo} de
comunicação (esconder padrões de tráfego), não à confidencialidade do
\textbf{conteúdo}.
\end{conceito}

\section{Gestão de riscos}

\begin{conceito}[Ameaça + vulnerabilidade = risco, resultando em impacto]
Um risco existe quando \textbf{uma ameaça encontra uma vulnerabilidade},
e disso resulta um \textbf{impacto} potencial. Retirar qualquer um dos
três elementos e o risco desaparece.

\begin{center}
\begin{tabular}{@{}p{3cm}p{9cm}@{}}
\toprule
\textbf{Ameaça} & o evento indesejado \textbf{em potencial} --- incêndio,
invasor, funcionário descuidado, \textit{ransomware}. Vem de fora do seu
controle direto: você \textbf{não elimina} a ameaça em si. \\
\textbf{Vulnerabilidade} & a fraqueza que permite a ameaça se
concretizar --- servidor sem \textit{patch}, senha fraca, ausência de
\textit{backup}. É o \textbf{único dos três} sobre o qual você age
diretamente. \\
\textbf{Impacto} & o dano, se a ameaça se concretizar. Você não impede o
impacto em si, mas pode \textbf{reduzi-lo} (segmentar a rede, cifrar o
dado que seria roubado) ou \textbf{transferi-lo} (seguro). \\
\bottomrule
\end{tabular}
\end{center}

O \textbf{risco} é a combinação de \textbf{probabilidade} $\times$
\textbf{impacto} --- é exatamente isso que uma matriz de risco 5$\times$5
desenha: probabilidade num eixo, impacto no outro, e a cor da célula
indica a prioridade de tratamento. Consequência importante: um evento
catastrófico mas muito improvável pode ficar \textbf{abaixo}, na
prioridade, de um evento pequeno mas frequente.
\end{conceito}

\begin{conceito}[As quatro formas de tratar o risco]
\begin{center}
\begin{tabular}{@{}p{2.6cm}p{9.4cm}@{}}
\toprule
\textbf{Mitigar} (reduzir) & implanta um controle e \textbf{diminui} a
probabilidade ou o impacto --- é o caso mais comum: aplicar
\textit{patch}, instalar um WAF, treinar a equipe. \\
\textbf{Transferir} (compartilhar) & passa a consequência
\textbf{financeira} a um terceiro --- seguro, terceirização com cláusula
de responsabilidade. Transferir \textbf{não reduz} a probabilidade de o
evento acontecer, só redistribui quem paga a conta. \\
\textbf{Evitar} & \textbf{deixa de fazer} a atividade que gera o risco ---
descontinuar um sistema legado, não coletar um dado sensível
desnecessário. \\
\textbf{Aceitar} (reter) & decide conviver com o risco,
\textbf{formalmente} e no nível hierárquico adequado, porque o controle
custaria mais que o dano potencial. \\
\bottomrule
\end{tabular}
\end{center}

\textbf{Risco residual} é o que \textbf{sobra depois} do tratamento
aplicado --- nunca é zero, e a norma exige que seja \textbf{aceito
formalmente} por quem tem alçada para isso; aceitar risco residual é
decisão de gestão, não do time técnico isoladamente.
\end{conceito}

\begin{cuidado}
Transferir o risco \textbf{não elimina} o risco --- o seguro paga o
prejuízo, mas o incidente acontece do mesmo jeito; quem reduz a
\textbf{chance} de acontecer é mitigar. \textbf{Aceitar} não é o mesmo
que \textbf{ignorar}: aceitar é um ato formal, documentado e aprovado;
deixar sem tratamento por simples omissão não é uma das quatro opções
formais.
\end{cuidado}

\section{Continuidade de negócio: RTO x RPO}

\begin{conceito}[Dois objetivos, guardados pela unidade do que se perde]
\begin{center}
\begin{tabular}{@{}p{1.6cm}p{4.4cm}p{6cm}@{}}
\toprule
\textbf{Sigla} & \textbf{Nome} & \textbf{O que limita} \\
\midrule
\textbf{RTO} & \textit{Recovery} \textbf{Time} \textit{Objective} &
\textbf{TEMPO}: quanto o serviço pode ficar indisponível até ser
restabelecido \\
\textbf{RPO} & \textit{Recovery} \textbf{Point} \textit{Objective} &
\textbf{DADO}: até que ponto no passado se aceita perder informação \\
\bottomrule
\end{tabular}
\end{center}

É o \textbf{RPO} que determina a \textbf{frequência do backup}: aceitar
perder no máximo 15 minutos de dado obriga a proteger os dados a cada 15
minutos, no máximo. O RTO, por sua vez, cobra da infraestrutura de
recuperação (redundância, sítio alternativo, tempo de restauração).

Siglas vizinhas frequentemente oferecidas junto: \textbf{MTBF}
(\textit{Mean Time Between Failures}, tempo médio \textbf{entre} falhas
sucessivas --- métrica de confiabilidade, não objetivo de plano);
\textbf{MTTR} (\textit{Mean Time To Repair}, tempo médio de reparo);
\textbf{SLA} (\textit{Service Level Agreement}, o \textbf{acordo}
contratual de nível de serviço, não uma métrica em si).
\end{conceito}

\begin{cuidado}
Ancore numa letra: \textbf{T} de RTO é \textbf{Tempo}; \textbf{P} de RPO
é \textbf{Ponto} no tempo (dado). Um cenário que dá dois números na ordem
"pode ficar fora X, pode perder Y" e depois inverte as siglas na resposta
é o erro clássico deste tema.
\end{cuidado}

\section{Detecção e resposta a incidentes}

\begin{conceito}
\textbf{IDS} (\textit{Intrusion Detection System}): \textbf{alerta} sobre
uma intrusão, de forma passiva --- não bloqueia. \textbf{IPS}
(\textit{Intrusion Prevention System}): \textbf{bloqueia} ativamente,
posicionado \textit{inline} no caminho do tráfego. \textbf{SIEM}
(\textit{Security Information and Event Management}): correlaciona
\textit{logs} e eventos de várias fontes para detecção e resposta.

\textbf{Resposta a incidentes} (referência: NIST SP 800-61), em quatro
fases ordenadas:
\begin{enumerate}
\item \textbf{Preparação}: antes de qualquer incidente --- time,
ferramentas, políticas prontas.
\item \textbf{Detecção e Análise}: identificar que um incidente está
ocorrendo e entender sua natureza.
\item \textbf{Contenção, Erradicação e Recuperação}: isolar o problema,
remover a causa raiz, restaurar serviços e dados.
\item \textbf{Atividade pós-incidente / Lições Aprendidas}: documentar o
ocorrido e ajustar controles para prevenir recorrência --- essa fase
realimenta a Preparação, fechando o ciclo.
\end{enumerate}
\end{conceito}

\begin{cuidado}
Detecção e Análise vêm \textbf{antes} da Contenção --- não se contém o
que ainda não foi identificado. Depois de conter e erradicar, vem
\textbf{recuperação seguida de lições aprendidas} --- essa ordem costuma
ser embaralhada em prova.
\end{cuidado}

\section{ISO/IEC 27001 e 27002}

\begin{conceito}[SGSI x guia de controles]
A \textbf{ISO/IEC 27001:2022} especifica os \textbf{requisitos} de um
Sistema de Gestão de Segurança da Informação (SGSI), incluindo seu Anexo
A. A \textbf{ISO/IEC 27002:2022} é o \textbf{guia de controles} detalhado
--- \textbf{93 controles}, organizados em \textbf{4 temas}:
organizacionais (37), pessoas (8), físicos (14), tecnológicos (34).

Exemplos de classificação temática: trabalho remoto $\to$ tema de
\textbf{pessoas}; mídia de armazenamento $\to$ tema \textbf{físico};
criptografia $\to$ tema \textbf{tecnológico}; definição de políticas
$\to$ tema \textbf{organizacional}.
\end{conceito}

\begin{conceito}[NIST Cybersecurity Framework e STRIDE]
\textbf{NIST CSF 1.1}: cinco funções --- \textit{Identify, Protect,
Detect, Respond, Recover}; a versão \textbf{2.0} acrescentou uma sexta
função, \textbf{Govern}, no centro do modelo.

\textbf{STRIDE} (modelagem de ameaças): seis categorias --- \textit{
Spoofing} (falsificar identidade), \textit{Tampering} (adulterar dado),
\textit{Repudiation} (negar uma ação realizada), \textit{Information
disclosure} (vazamento de informação), \textit{Denial of service}
(indisponibilidade forçada), \textit{Elevation of privilege} (ganhar
mais permissão do que deveria).
\end{conceito}
